\begin{titlepage}
\centering
{\Large 技术综述教学笔记\par}
\vspace{1.0cm}
{\huge\bfseries 开源一段论文探索之旅: AI 变迁史的四条主线}\par
\vspace{0.5cm}
{\Large Zhang Xiaojun 对话谢青池\par}
\vspace{0.35cm}
{\large 基于 Zhang Xiaojun Podcast 公开视频、GPU 转写稿、视频描述论文路线图与自制概念图整理\par}
\vspace{0.3cm}
{\large 2025-10-28\par}
\vspace{0.85cm}
\includegraphics[width=0.88\textwidth,height=0.42\textheight,keepaspectratio]{cover.jpg}\par
\vfill
\begin{tcolorbox}[width=0.92\textwidth, colback=black!2!white, colframe=black!55, sharp corners]
\textbf{视频标题}: 117. 开源一段论文探索之旅: 模型范式、Infra和数据、语言、多模态的完整变迁史\par
\textbf{视频频道}: Zhang Xiaojun Podcast\par
\textbf{Duración del video}: 04:22:37\par
\textbf{Enlace del video}: \href{\videourl}{\nolinkurl{\videourl}}\par
\textbf{素材说明}: YouTube 无可用字幕；正文基于 faster-whisper large-v3 CUDA 转写、YouTube 描述中的 36 篇论文路线图、章节结构、封面和自制教学图整理.YouTube 版为静态封面/固定视觉, 不重复插入人物帧；描述中的 Bilibili 投屏版在本机抓取元数据返回 HTTP 412, 飞书 PPT 抽取未拿到正文, 故当前版本以可靠本地源为准.\par
\textbf{Página del proyecto}: \href{\repourl}{\nolinkurl{\repourl}}\par
\end{tcolorbox}
\end{titlepage}

\tableofcontents
\newpage

